\documentclass[a4paper,11pt]{article}

\usepackage[utf8]{inputenc}
\usepackage[T1]{fontenc}
\usepackage{helvet}
\renewcommand{\familydefault}{\sfdefault}
\usepackage[a4paper, top=2.5cm, bottom=2.5cm, left=2.5cm, right=2.5cm]{geometry}
\usepackage{fancyhdr}
\usepackage{booktabs}
\usepackage{tabularx}
\usepackage{longtable}
\usepackage{array}
\usepackage{xcolor}
\usepackage{enumitem}
\usepackage{hyperref}
\usepackage{pifont}
\usepackage{amssymb}
\usepackage{graphicx}
\usepackage{multirow}
\usepackage{colortbl}

\hypersetup{
	colorlinks=true,
	linkcolor=black,
	urlcolor=blue
}

% Custom colors
\definecolor{headerblue}{RGB}{0,70,127}
\definecolor{lightgray}{RGB}{240,240,240}
\definecolor{passgreen}{RGB}{0,150,0}
\definecolor{failred}{RGB}{180,0,0}
\definecolor{sectionblue}{RGB}{220,235,250}

% Custom commands
\newcommand{\testpass}{\textcolor{passgreen}{\ding{51}~Pass}}
\newcommand{\testfail}{\textcolor{failred}{\ding{55}~Fail}}
\newcommand{\checkbox}{$\square$}
\newcommand{\key}[1]{\fbox{\texttt{#1}}}
\newcommand{\button}[1]{\textbf{``#1''}}
\newcommand{\menu}[1]{\textit{#1}}
\newcommand{\tab}[1]{\textbf{#1 tab}}
\newcommand{\ui}[1]{\textsf{#1}}

% Test case IDs are derived automatically from the section/subsection number,
% e.g. the second test in section 3 prints as "TC-3.2". Adding or removing a test
% only renumbers the tests inside that same section; the summary table and all
% cross-references update on the next compile via \ref{tc:<key>}.
\makeatletter
\newcommand{\testcase}[2]{% #1 = label key, #2 = test title
    \refstepcounter{subsection}%
    \protected@edef\@currentlabel{TC-\thesubsection}%
    \subsection*{TC-\thesubsection: #2}%
    \addcontentsline{toc}{subsection}{TC-\thesubsection: #2}%
    \label{tc:#1}%
}
\makeatother

% Page style
\pagestyle{fancy}
\fancyhf{}
\fancyhead[L]{\small PRo3D Test Protocol}
\fancyhead[R]{\small Version \testversion}
\fancyfoot[L]{\small \today}
\fancyfoot[R]{\small Page \thepage}
\renewcommand{\headrulewidth}{0.4pt}
\renewcommand{\footrulewidth}{0.4pt}

\newcommand{\testversion}{1.0}

%----------------------------------------------------------------------------------------
\begin{document}
	%----------------------------------------------------------------------------------------
	
	%----------------------------------------------------------------------------------------
	% TITLE PAGE
	%----------------------------------------------------------------------------------------
	\begin{titlepage}
		\centering
		\vspace*{2cm}
		{\scshape\Huge \textcolor{headerblue}{\textbf{PRo3D}} \par}
		\vspace{0.5cm}
		{\scshape\LARGE Test Protocol \par}
		\vspace{0.3cm}
		{\large Version \testversion \par}
		\vspace{2cm}
		\rule{\textwidth}{0.5pt}\\[0.4cm]
		{\large\textbf{Purpose:} Manual feature verification against the PRo3D Short User Manual \par}
		\rule{\textwidth}{0.5pt}\\[1cm]
		
		\begin{tabular}{ll}
			\textbf{Document Status:} & Draft \\[4pt]
			\textbf{Created:}         & \today \\[4pt]
			\textbf{Reference Manual:} & PRo3D Short User Manual v5.5.0 \\[4pt]
			\textbf{Scope:}           & All features described in the Short User Manual \\
		\end{tabular}
		
		\vspace{2cm}
		\textbf{Tester Sign-off}\\[1em]
		\begin{tabular}{p{5cm} p{5cm} p{3cm}}
			\hline\\[-6pt]
			Name & Signature & Date \\[12pt]
			\hline\\
			& & \\[12pt]
			\hline
		\end{tabular}
		
		\vfill
		{\small Planetary Robotics 3D Viewer $\cdot$ VRVis Research Center}
	\end{titlepage}
	
	\newpage
	
	%----------------------------------------------------------------------------------------
	% HOW TO USE THIS DOCUMENT
	%----------------------------------------------------------------------------------------
	\section*{How to Use This Test Protocol}
	
	This document provides step-by-step test cases covering all features described in the PRo3D Short User Manual.
	Each test case specifies:
	\begin{itemize}
		\item \textbf{Prerequisites} -- what must be set up before starting the test.
		\item \textbf{Steps} -- the exact sequence of actions (buttons to click, keys to press).
		\item \textbf{Expected Result} -- what the tester should observe if the feature works correctly.
		\item \textbf{Result} -- a checkbox the tester fills in as \testpass{} or \testfail{}.
		\item \textbf{Notes} -- space to record defects, screenshots, or observations.
	\end{itemize}
	
	\noindent\textbf{Notation used in this document:}
	\begin{itemize}
		\item \button{Button Name} -- a button in the GUI that the tester should click.
		\item \menu{Menu $\rightarrow$ Item} -- a menu path to follow.
		\item \key{KEY} -- a keyboard key to press.
		\item \key{CTRL}+\key{LMB} -- hold Ctrl and click the Left Mouse Button on the 3D surface.
		\item \tab{Tab Name} -- a tab in the right-hand side panel.
	\end{itemize}
	
	\noindent A test case \textbf{passes} when all expected results are observed. If any expected result is not observed, mark the test as \textbf{fail} and add details in the Notes field.
	
	\newpage
	
	%----------------------------------------------------------------------------------------
	% SUMMARY TABLE
	%----------------------------------------------------------------------------------------
	\section*{Test Summary}
	\label{sec:summary}
	
	\begin{longtable}{p{1.8cm} p{6.5cm} p{2.5cm} p{2.5cm}}
		\toprule
		\textbf{Test ID} & \textbf{Feature} & \textbf{Result} & \textbf{Notes} \\
		\midrule
		\endhead
		\ref{tc:01} & Start PRo3D & \checkbox~Pass / \checkbox~Fail & \\
		\ref{tc:02} & Add Surface (OPC) & \checkbox~Pass / \checkbox~Fail & \\
		\ref{tc:03} & Surface FlyTo & \checkbox~Pass / \checkbox~Fail & \\
		\ref{tc:04} & Save Scene & \checkbox~Pass / \checkbox~Fail & \\
		\ref{tc:05} & Load Scene (Open) & \checkbox~Pass / \checkbox~Fail & \\
		\ref{tc:06} & Load Scene (Recent) & \checkbox~Pass / \checkbox~Fail & \\
		\ref{tc:07} & FreeFly Navigation & \checkbox~Pass / \checkbox~Fail & \\
		\ref{tc:08} & ArcBall Navigation & \checkbox~Pass / \checkbox~Fail & \\
		\ref{tc:09} & Pick Explore Center & \checkbox~Pass / \checkbox~Fail & \\
		\ref{tc:10} & Place Coordinate System & \checkbox~Pass / \checkbox~Fail & \\
		\ref{tc:11} & Draw Point Annotation & \checkbox~Pass / \checkbox~Fail & \\
		\ref{tc:12} & Draw Line Annotation & \checkbox~Pass / \checkbox~Fail & \\
		\ref{tc:13} & Draw Polyline Annotation & \checkbox~Pass / \checkbox~Fail & \\
		\ref{tc:14} & Draw Polygon Annotation & \checkbox~Pass / \checkbox~Fail & \\
		\ref{tc:15} & Draw Dip and Strike Annotation & \checkbox~Pass / \checkbox~Fail & \\
		\ref{tc:16} & Annotation Projection Modes & \checkbox~Pass / \checkbox~Fail & \\
		\ref{tc:17} & Pick Annotation (via UI list) & \checkbox~Pass / \checkbox~Fail & \\
		\ref{tc:18} & Pick Annotation (via 3D view) & \checkbox~Pass / \checkbox~Fail & \\
		\ref{tc:19} & Pick Surface (via UI list) & \checkbox~Pass / \checkbox~Fail & \\
		\ref{tc:20} & Pick Surface (via 3D view) & \checkbox~Pass / \checkbox~Fail & \\
		\ref{tc:21} & Surface Properties Panel & \checkbox~Pass / \checkbox~Fail & \\
		\ref{tc:22} & Surface Visibility Toggle & \checkbox~Pass / \checkbox~Fail & \\
		\ref{tc:23} & Surface Color Correction & \checkbox~Pass / \checkbox~Fail & \\
		\ref{tc:24} & Surface Translation & \checkbox~Pass / \checkbox~Fail & \\
		\ref{tc:25} & Surface FillMode (Wireframe) & \checkbox~Pass / \checkbox~Fail & \\
		\ref{tc:26} & Annotation Properties Panel & \checkbox~Pass / \checkbox~Fail & \\
		\ref{tc:27} & Annotation Measurements & \checkbox~Pass / \checkbox~Fail & \\
		\ref{tc:28} & Annotation Text Note & \checkbox~Pass / \checkbox~Fail & \\
		\ref{tc:29} & Place Scalebar & \checkbox~Pass / \checkbox~Fail & \\
		\ref{tc:30} & Scalebar Properties & \checkbox~Pass / \checkbox~Fail & \\
		\ref{tc:31} & Add Bookmark & \checkbox~Pass / \checkbox~Fail & \\
		\ref{tc:32} & Bookmark FlyTo & \checkbox~Pass / \checkbox~Fail & \\
		\ref{tc:33} & Sequenced Bookmarks -- Add \& Animate & \checkbox~Pass / \checkbox~Fail & \\
		\ref{tc:34} & Sequenced Bookmarks -- Record \& Render & \checkbox~Pass / \checkbox~Fail & \\
		\ref{tc:35} & Viewer Config Settings & \checkbox~Pass / \checkbox~Fail & \\
		\ref{tc:36} & Screenshots & \checkbox~Pass / \checkbox~Fail & \\
		\ref{tc:37} & Grouping -- Create Group & \checkbox~Pass / \checkbox~Fail & \\
		\ref{tc:38} & Grouping -- Move Leafs & \checkbox~Pass / \checkbox~Fail & \\
		\ref{tc:39} & View Planner -- Place Rover & \checkbox~Pass / \checkbox~Fail & \\
		\ref{tc:40} & View Planner -- Instrument View & \checkbox~Pass / \checkbox~Fail & \\
		\ref{tc:41} & GIS View -- Load SPICE Kernel & \checkbox~Pass / \checkbox~Fail & \\
		\ref{tc:42} & GIS View -- Observation Settings & \checkbox~Pass / \checkbox~Fail & \\
		\ref{tc:43} & Contour Lines & \checkbox~Pass / \checkbox~Fail & \\
		\ref{tc:44} & Multitexturing & \checkbox~Pass / \checkbox~Fail & \\
		\ref{tc:57} & Cross Sections & \checkbox~Pass / \checkbox~Fail & \\
		\ref{tc:45} & Surface Comparison -- Activate Mode & \checkbox~Pass / \checkbox~Fail & \\
		\ref{tc:46} & Surface Comparison -- Select Area & \checkbox~Pass / \checkbox~Fail & \\
		\ref{tc:47} & Surface Comparison -- Surface Measurements & \checkbox~Pass / \checkbox~Fail & \\
		\ref{tc:48} & Surface Comparison -- Length Measurements & \checkbox~Pass / \checkbox~Fail & \\
		\ref{tc:49} & Minerva -- Load Data \& Listing & \checkbox~Pass / \checkbox~Fail & \\
		\ref{tc:50} & Minerva -- Query App Filter & \checkbox~Pass / \checkbox~Fail & \\
		\ref{tc:51} & Minerva -- Product Selection & \checkbox~Pass / \checkbox~Fail & \\
		\ref{tc:52} & Minerva -- Pick Linking & \checkbox~Pass / \checkbox~Fail & \\
		\ref{tc:53} & CLI -- Snapshot Rendering & \checkbox~Pass / \checkbox~Fail & \\
		\ref{tc:54} & RIMFAX -- Load Traverse & \checkbox~Pass / \checkbox~Fail & \\
		\ref{tc:55} & RIMFAX -- Load Surfaces & \checkbox~Pass / \checkbox~Fail & \\
		\ref{tc:56} & Keyboard Shortcuts & \checkbox~Pass / \checkbox~Fail & \\
		\bottomrule
	\end{longtable}
	
	\newpage
	
	%----------------------------------------------------------------------------------------
	% TABLE OF CONTENTS
	%----------------------------------------------------------------------------------------
	\tableofcontents
	\newpage
	
	%----------------------------------------------------------------------------------------
	%  SECTION 1: STARTING PRo3D
	%----------------------------------------------------------------------------------------
	\section{Starting PRo3D}
	
	%--- tc:01 ---
	\testcase{01}{Start PRo3D}
	
	\begin{tabularx}{\textwidth}{lX}
		\toprule
		\textbf{Prerequisites} & PRo3D is installed. \\
		\midrule
		\textbf{Steps} &
		\begin{enumerate}[leftmargin=*, nosep]
			\item Double-click \texttt{PRo3D.exe} (or the application shortcut).
			\item Wait for the application to finish loading.
		\end{enumerate} \\
		\midrule
		\textbf{Expected Result} &
		\begin{itemize}[leftmargin=*, nosep]
			\item The PRo3D main window opens.
			\item A menu icon (hamburger or similar) is visible in the top-left corner of the window.
			\item No error messages or crash dialogs appear.
		\end{itemize} \\
		\midrule
		\textbf{Result} & \checkbox~Pass \quad \checkbox~Fail \\
		\midrule
		\textbf{Notes} & \\[1.5cm]
		\bottomrule
	\end{tabularx}
	
	\vspace{1em}
	
	%--- tc:02 ---
	\testcase{02}{Add Surface (OPC)}
	
	\begin{tabularx}{\textwidth}{lX}
		\toprule
		\textbf{Prerequisites} & PRo3D is running (\ref{tc:01} passed). An OPC surface folder is available on disk. \\
		\midrule
		\textbf{Steps} &
		\begin{enumerate}[leftmargin=*, nosep]
			\item Click the main menu icon in the top-left corner.
			\item Hover over the \menu{Import} tab.
			\item Click \button{OPC}.
			\item In the ``Select Folder'' dialog, navigate to the OPC folder and click \button{Select Folder}.
			\item Wait for the surface to load.
		\end{enumerate} \\
		\midrule
		\textbf{Expected Result} &
		\begin{itemize}[leftmargin=*, nosep]
			\item The ``Select Folder'' dialog opens and allows folder selection.
			\item After confirming, the surface is listed in the \tab{Surfaces} panel on the right side of the window.
			\item The surface name appears in the list.
		\end{itemize} \\
		\midrule
		\textbf{Result} & \checkbox~Pass \quad \checkbox~Fail \\
		\midrule
		\textbf{Notes} & \\[1.5cm]
		\bottomrule
	\end{tabularx}
	
	\vspace{1em}
	
	%--- tc:03 ---
	\testcase{03}{Surface FlyTo}
	
	\begin{tabularx}{\textwidth}{lX}
		\toprule
		\textbf{Prerequisites} & A surface has been loaded (\ref{tc:02} passed). \\
		\midrule
		\textbf{Steps} &
		\begin{enumerate}[leftmargin=*, nosep]
			\item In the \tab{Surfaces} panel (right side), locate the loaded surface in the list.
			\item Click the \button{FlyTo} button below the surface name.
		\end{enumerate} \\
		\midrule
		\textbf{Expected Result} &
		\begin{itemize}[leftmargin=*, nosep]
			\item The camera performs a smooth fly-to animation.
			\item The 3D surface becomes visible and centered in the main 3D view.
		\end{itemize} \\
		\midrule
		\textbf{Result} & \checkbox~Pass \quad \checkbox~Fail \\
		\midrule
		\textbf{Notes} & \\[1.5cm]
		\bottomrule
	\end{tabularx}
	
	\vspace{1em}
	
	%--- tc:04 ---
	\testcase{04}{Save Scene}
	
	\begin{tabularx}{\textwidth}{lX}
		\toprule
		\textbf{Prerequisites} & At least one surface is loaded (\ref{tc:02} passed). \\
		\midrule
		\textbf{Steps} &
		\begin{enumerate}[leftmargin=*, nosep]
			\item Click the main menu icon in the top-left corner.
			\item Navigate to \menu{Scene $\rightarrow$ Save as\ldots}.
			\item In the Save dialog, enter a name for the scene file (e.g., \texttt{TestScene}).
			\item Click \button{Save}.
		\end{enumerate} \\
		\midrule
		\textbf{Expected Result} &
		\begin{itemize}[leftmargin=*, nosep]
			\item The save dialog opens.
			\item After saving, a \texttt{.scn} file appears in the chosen directory.
			\item No error messages are shown.
		\end{itemize} \\
		\midrule
		\textbf{Result} & \checkbox~Pass \quad \checkbox~Fail \\
		\midrule
		\textbf{Notes} & \\[1.5cm]
		\bottomrule
	\end{tabularx}
	
	\vspace{1em}
	
	%--- tc:05 ---
	\testcase{05}{Load Scene (Open)}
	
	\begin{tabularx}{\textwidth}{lX}
		\toprule
		\textbf{Prerequisites} & A scene file (\texttt{.scn}) exists on disk (\ref{tc:04} passed). \\
		\midrule
		\textbf{Steps} &
		\begin{enumerate}[leftmargin=*, nosep]
			\item Restart PRo3D (or close the current scene).
			\item Click the main menu icon in the top-left.
			\item Navigate to \menu{Scene $\rightarrow$ Open}.
			\item In the file dialog, navigate to the saved \texttt{.scn} file and select it.
			\item Click \button{Open} (or double-click the file).
		\end{enumerate} \\
		\midrule
		\textbf{Expected Result} &
		\begin{itemize}[leftmargin=*, nosep]
			\item The scene loads successfully.
			\item Previously added surfaces appear in the \tab{Surfaces} list.
			\item The 3D view displays the loaded surface data.
		\end{itemize} \\
		\midrule
		\textbf{Result} & \checkbox~Pass \quad \checkbox~Fail \\
		\midrule
		\textbf{Notes} & \\[1.5cm]
		\bottomrule
	\end{tabularx}
	
	\vspace{1em}
	
	%--- tc:06 ---
	\testcase{06}{Load Scene (Recent)}
	
	\begin{tabularx}{\textwidth}{lX}
		\toprule
		\textbf{Prerequisites} & A scene has been opened at least once before (\ref{tc:05} passed). \\
		\midrule
		\textbf{Steps} &
		\begin{enumerate}[leftmargin=*, nosep]
			\item Click the main menu icon in the top-left.
			\item Hover over \menu{Scene $\rightarrow$ Recent}.
			\item A list of recently opened scenes appears; click the desired scene name.
		\end{enumerate} \\
		\midrule
		\textbf{Expected Result} &
		\begin{itemize}[leftmargin=*, nosep]
			\item A list of recent scene files is shown on hovering.
			\item Clicking a scene name loads that scene without opening a file dialog.
			\item Surfaces from the selected scene appear in the list and 3D view.
		\end{itemize} \\
		\midrule
		\textbf{Result} & \checkbox~Pass \quad \checkbox~Fail \\
		\midrule
		\textbf{Notes} & \\[1.5cm]
		\bottomrule
	\end{tabularx}
	
	%----------------------------------------------------------------------------------------
	% SECTION 2: VIEWER ACTIONS / NAVIGATION
	%----------------------------------------------------------------------------------------
	\newpage
	\section{Viewer Actions and Navigation}
	
	%--- tc:07 ---
	\testcase{07}{FreeFly Navigation}
	
	\begin{tabularx}{\textwidth}{lX}
		\toprule
		\textbf{Prerequisites} & A surface is loaded and visible in the 3D view (\ref{tc:03} passed). The navigation mode is set to FreeFly. \\
		\midrule
		\textbf{Steps} &
		\begin{enumerate}[leftmargin=*, nosep]
			\item Ensure the viewer interaction mode is set to \button{PickExploreCenter} (press \key{F1} or select from the actions menu).
			\item Press \key{W} to move forward.
			\item Press \key{S} to move backward.
			\item Press \key{A} to strafe left.
			\item Press \key{D} to strafe right.
			\item Hold \key{LMB} (Left Mouse Button) and drag to rotate the camera.
			\item Scroll the mouse wheel to zoom in and out.
			\item Hold \key{RMB} (Right Mouse Button) and drag to move forward/backward.
			\item Hold the Middle Mouse Button and drag to pan.
		\end{enumerate} \\
		\midrule
		\textbf{Expected Result} &
		\begin{itemize}[leftmargin=*, nosep]
			\item \key{W}/\key{S} moves the camera forward and backward through the scene.
			\item \key{A}/\key{D} strafes the camera left and right.
			\item Dragging with \key{LMB} held rotates the camera's orientation.
			\item Mouse wheel zooms smoothly in and out.
			\item Dragging with \key{RMB} held moves the camera forward/backward.
			\item Dragging with Middle Mouse Button pans the view.
		\end{itemize} \\
		\midrule
		\textbf{Result} & \checkbox~Pass \quad \checkbox~Fail \\
		\midrule
		\textbf{Notes} & \\[1.5cm]
		\bottomrule
	\end{tabularx}
	
	\vspace{1em}
	
	%--- tc:08 ---
	\testcase{08}{ArcBall Navigation}
	
	\begin{tabularx}{\textwidth}{lX}
		\toprule
		\textbf{Prerequisites} & ArcBall navigation mode is active. A surface is loaded. \\
		\midrule
		\textbf{Steps} &
		\begin{enumerate}[leftmargin=*, nosep]
			\item Switch to ArcBall mode via the navigation mode selector.
			\item Hold \key{LMB} and drag to rotate the camera around the current explore center.
			\item Press \key{W} or scroll the mouse wheel to move forward.
			\item Press \key{S} or scroll the mouse wheel to move backward.
		\end{enumerate} \\
		\midrule
		\textbf{Expected Result} &
		\begin{itemize}[leftmargin=*, nosep]
			\item The camera orbits around the explore center when \key{LMB} is held and dragged.
			\item \key{W}/\key{S} and the mouse wheel change the distance to the explore center.
			\item The explore center remains stationary during rotation.
		\end{itemize} \\
		\midrule
		\textbf{Result} & \checkbox~Pass \quad \checkbox~Fail \\
		\midrule
		\textbf{Notes} & \\[1.5cm]
		\bottomrule
	\end{tabularx}
	
	\vspace{1em}
	
	%--- tc:09 ---
	\testcase{09}{Pick Explore Center}
	
	\begin{tabularx}{\textwidth}{lX}
		\toprule
		\textbf{Prerequisites} & ArcBall navigation mode is active. A surface is loaded and visible. \\
		\midrule
		\textbf{Steps} &
		\begin{enumerate}[leftmargin=*, nosep]
			\item In the actions drop-down, select \button{PickExploreCenter} (or press \key{F1}).
			\item Hold \key{Ctrl} and click \key{LMB} on a point on the 3D surface.
			\item Observe the viewport.
			\item Then hold \key{LMB} and drag to orbit.
		\end{enumerate} \\
		\midrule
		\textbf{Expected Result} &
		\begin{itemize}[leftmargin=*, nosep]
			\item A pink dot appears at the clicked location on the surface.
			\item Orbiting (dragging with \key{LMB}) now rotates the camera around the new pink dot position.
		\end{itemize} \\
		\midrule
		\textbf{Result} & \checkbox~Pass \quad \checkbox~Fail \\
		\midrule
		\textbf{Notes} & \\[1.5cm]
		\bottomrule
	\end{tabularx}
	
	\vspace{1em}
	
	%--- tc:10 ---
	\testcase{10}{Place Coordinate System}
	
	\begin{tabularx}{\textwidth}{lX}
		\toprule
		\textbf{Prerequisites} & A surface is loaded and visible. \\
		\midrule
		\textbf{Steps} &
		\begin{enumerate}[leftmargin=*, nosep]
			\item In the actions drop-down, select \button{PlaceCoordinateSystem} (or press \key{F4}).
			\item Select a unit of measurement from the unit drop-down (e.g., meters).
			\item Hold \key{Ctrl} and click \key{LMB} on a point on the 3D surface.
			\item Observe the viewport and the Config panel.
		\end{enumerate} \\
		\midrule
		\textbf{Expected Result} &
		\begin{itemize}[leftmargin=*, nosep]
			\item A coordinate system widget appears at the clicked position on the surface.
			\item A blue arrow (up vector) points in the positive Z direction.
			\item A red arrow (north vector) points in the positive Y direction.
			\item The position, up vector, and north vector values are updated in the \tab{Config} panel under \menu{Coordinate System}.
		\end{itemize} \\
		\midrule
		\textbf{Result} & \checkbox~Pass \quad \checkbox~Fail \\
		\midrule
		\textbf{Notes} & \\[1.5cm]
		\bottomrule
	\end{tabularx}
	
	%----------------------------------------------------------------------------------------
	% SECTION 3: ANNOTATIONS
	%----------------------------------------------------------------------------------------
	\newpage
	\section{Drawing and Managing Annotations}
	
	%--- tc:11 ---
	\testcase{11}{Draw Point Annotation}
	
	\begin{tabularx}{\textwidth}{lX}
		\toprule
		\textbf{Prerequisites} & A surface is loaded and visible. \\
		\midrule
		\textbf{Steps} &
		\begin{enumerate}[leftmargin=*, nosep]
			\item In the actions drop-down, select \button{DrawAnnotation} (or press \key{F2}).
			\item In the annotation mode selector, choose \button{Point}.
			\item Hold \key{Ctrl} and click \key{LMB} on the surface.
		\end{enumerate} \\
		\midrule
		\textbf{Expected Result} &
		\begin{itemize}[leftmargin=*, nosep]
			\item A single point marker appears on the surface at the clicked location.
			\item The new annotation appears in the \tab{Annotations} list on the right.
			\item Clicking the annotation name in the list highlights it in green.
		\end{itemize} \\
		\midrule
		\textbf{Result} & \checkbox~Pass \quad \checkbox~Fail \\
		\midrule
		\textbf{Notes} & \\[1.5cm]
		\bottomrule
	\end{tabularx}
	
	\vspace{1em}
	
	%--- tc:12 ---
	\testcase{12}{Draw Line Annotation}
	
	\begin{tabularx}{\textwidth}{lX}
		\toprule
		\textbf{Prerequisites} & DrawAnnotation mode is active. \\
		\midrule
		\textbf{Steps} &
		\begin{enumerate}[leftmargin=*, nosep]
			\item In the annotation mode selector, choose \button{Line}.
			\item Hold \key{Ctrl} and click \key{LMB} on a first point on the surface.
			\item Hold \key{Ctrl} and click \key{LMB} on a second point on the surface.
		\end{enumerate} \\
		\midrule
		\textbf{Expected Result} &
		\begin{itemize}[leftmargin=*, nosep]
			\item Two points are placed and connected by a line projected onto the surface.
			\item The line appears in the \tab{Annotations} list.
			\item The \menu{Measurements} panel shows distance values (Length, WayLength).
		\end{itemize} \\
		\midrule
		\textbf{Result} & \checkbox~Pass \quad \checkbox~Fail \\
		\midrule
		\textbf{Notes} & \\[1.5cm]
		\bottomrule
	\end{tabularx}
	
	\vspace{1em}
	
	%--- tc:13 ---
	\testcase{13}{Draw Polyline Annotation}
	
	\begin{tabularx}{\textwidth}{lX}
		\toprule
		\textbf{Prerequisites} & DrawAnnotation mode is active. \\
		\midrule
		\textbf{Steps} &
		\begin{enumerate}[leftmargin=*, nosep]
			\item In the annotation mode selector, choose \button{Polyline}.
			\item Hold \key{Ctrl} and click \key{LMB} to place the first point.
			\item Hold \key{Ctrl} and click \key{LMB} to place at least two more points.
			\item Press \key{Enter} to finish the polyline.
		\end{enumerate} \\
		\midrule
		\textbf{Expected Result} &
		\begin{itemize}[leftmargin=*, nosep]
			\item Each click places a new point and a connecting line segment is drawn.
			\item Pressing \key{Enter} finalizes the polyline.
			\item The polyline appears in the \tab{Annotations} list.
			\item Measurements (Length, WayLength, Slope, etc.) are shown in the properties panel.
		\end{itemize} \\
		\midrule
		\textbf{Result} & \checkbox~Pass \quad \checkbox~Fail \\
		\midrule
		\textbf{Notes} & \\[1.5cm]
		\bottomrule
	\end{tabularx}
	
	\vspace{1em}
	
	%--- tc:14 ---
	\testcase{14}{Draw Polygon Annotation}
	
	\begin{tabularx}{\textwidth}{lX}
		\toprule
		\textbf{Prerequisites} & DrawAnnotation mode is active. \\
		\midrule
		\textbf{Steps} &
		\begin{enumerate}[leftmargin=*, nosep]
			\item In the annotation mode selector, choose \button{Polygon}.
			\item Hold \key{Ctrl} and click \key{LMB} to place at least three points.
			\item Press \key{Enter} to close and finish the polygon.
		\end{enumerate} \\
		\midrule
		\textbf{Expected Result} &
		\begin{itemize}[leftmargin=*, nosep]
			\item Points are placed and connected by segments.
			\item Pressing \key{Enter} closes the polygon (last point connects to first point).
			\item The polygon appears in the \tab{Annotations} list.
		\end{itemize} \\
		\midrule
		\textbf{Result} & \checkbox~Pass \quad \checkbox~Fail \\
		\midrule
		\textbf{Notes} & \\[1.5cm]
		\bottomrule
	\end{tabularx}
	
	\vspace{1em}
	
	%--- tc:15 ---
	\testcase{15}{Draw Dip and Strike (DnS) Annotation}
	
	\begin{tabularx}{\textwidth}{lX}
		\toprule
		\textbf{Prerequisites} & DrawAnnotation mode is active. A surface with geological structure is loaded. \\
		\midrule
		\textbf{Steps} &
		\begin{enumerate}[leftmargin=*, nosep]
			\item In the annotation mode selector, choose \button{DnS}.
			\item Hold \key{Ctrl} and click \key{LMB} to place at least three points along a rock layer.
			\item Press \key{Enter} to finish and compute the dip and strike.
		\end{enumerate} \\
		\midrule
		\textbf{Expected Result} &
		\begin{itemize}[leftmargin=*, nosep]
			\item A polyline (blue) is drawn along the picked points.
			\item A fitted plane is shown with a strike vector (red) and dip vector (green).
			\item A disc visualizing the dip direction appears on the surface.
			\item Dip and Strike measurements (azimuth, dip angle) appear in the properties panel.
			\item The annotation appears in the \tab{Annotations} list.
		\end{itemize} \\
		\midrule
		\textbf{Result} & \checkbox~Pass \quad \checkbox~Fail \\
		\midrule
		\textbf{Notes} & \\[1.5cm]
		\bottomrule
	\end{tabularx}
	
	\vspace{1em}
	
	%--- tc:16 ---
	\testcase{16}{Annotation Projection Modes}
	
	\begin{tabularx}{\textwidth}{lX}
		\toprule
		\textbf{Prerequisites} & DrawAnnotation mode is active in Line or Polyline mode. \\
		\midrule
		\textbf{Steps} &
		\begin{enumerate}[leftmargin=*, nosep]
			\item In the projection selector, choose \button{Linear}.
			\item Draw a line annotation on the surface. Observe the result.
			\item Draw another line with \button{ViewPoint} projection and compare.
			\item Draw another line with \button{Sky} projection and compare.
		\end{enumerate} \\
		\midrule
		\textbf{Expected Result} &
		\begin{itemize}[leftmargin=*, nosep]
			\item \textbf{Linear}: The line is a straight point-to-point connection without surface projection.
			\item \textbf{ViewPoint}: The line follows the surface contour as seen from the current camera position; it hugs depressions and ridges.
			\item \textbf{Sky}: The line is projected along the scene's up-vector, following the surface from above.
			\item The three lines visibly differ in their path across the surface.
		\end{itemize} \\
		\midrule
		\textbf{Result} & \checkbox~Pass \quad \checkbox~Fail \\
		\midrule
		\textbf{Notes} & \\[1.5cm]
		\bottomrule
	\end{tabularx}
	
	\vspace{1em}
	
	%--- tc:17 ---
	\testcase{17}{Pick Annotation (via UI List)}
	
	\begin{tabularx}{\textwidth}{lX}
		\toprule
		\textbf{Prerequisites} & At least one annotation exists (\ref{tc:11} through \ref{tc:15} passed). \\
		\midrule
		\textbf{Steps} &
		\begin{enumerate}[leftmargin=*, nosep]
			\item In the \tab{Annotations} panel, click on an annotation's name.
		\end{enumerate} \\
		\midrule
		\textbf{Expected Result} &
		\begin{itemize}[leftmargin=*, nosep]
			\item The annotation's name turns green in the list.
			\item The annotation is highlighted with a green border in the 3D view.
			\item The annotation's properties appear in the Properties panel.
		\end{itemize} \\
		\midrule
		\textbf{Result} & \checkbox~Pass \quad \checkbox~Fail \\
		\midrule
		\textbf{Notes} & \\[1.5cm]
		\bottomrule
	\end{tabularx}
	
	\vspace{1em}
	
	%--- tc:18 ---
	\testcase{18}{Pick Annotation (via 3D View)}
	
	\begin{tabularx}{\textwidth}{lX}
		\toprule
		\textbf{Prerequisites} & At least one annotation exists in the 3D view. \\
		\midrule
		\textbf{Steps} &
		\begin{enumerate}[leftmargin=*, nosep]
			\item In the actions drop-down, select \button{PickAnnotation} (or press \key{F3}).
			\item Hold \key{Ctrl} and click \key{LMB} on or near an annotation in the main 3D view.
		\end{enumerate} \\
		\midrule
		\textbf{Expected Result} &
		\begin{itemize}[leftmargin=*, nosep]
			\item The annotation under the cursor is selected.
			\item It turns green in the annotations list.
			\item It gets a green border in the 3D view.
			\item Its properties are shown in the Properties panel.
		\end{itemize} \\
		\midrule
		\textbf{Result} & \checkbox~Pass \quad \checkbox~Fail \\
		\midrule
		\textbf{Notes} & \\[1.5cm]
		\bottomrule
	\end{tabularx}
	
	\vspace{1em}
	
	%--- tc:19 ---
	\testcase{19}{Pick Surface (via UI List)}
	
	\begin{tabularx}{\textwidth}{lX}
		\toprule
		\textbf{Prerequisites} & At least one surface is loaded. \\
		\midrule
		\textbf{Steps} &
		\begin{enumerate}[leftmargin=*, nosep]
			\item Navigate to the \tab{Surfaces} panel.
			\item Click on the surface's name in the list.
		\end{enumerate} \\
		\midrule
		\textbf{Expected Result} &
		\begin{itemize}[leftmargin=*, nosep]
			\item The surface's name turns green in the list.
			\item The surface's properties appear in the Properties panel.
		\end{itemize} \\
		\midrule
		\textbf{Result} & \checkbox~Pass \quad \checkbox~Fail \\
		\midrule
		\textbf{Notes} & \\[1.5cm]
		\bottomrule
	\end{tabularx}
	
	\vspace{1em}
	
	%--- tc:20 ---
	\testcase{20}{Pick Surface (via 3D View)}
	
	\begin{tabularx}{\textwidth}{lX}
		\toprule
		\textbf{Prerequisites} & At least one surface is loaded and visible. \\
		\midrule
		\textbf{Steps} &
		\begin{enumerate}[leftmargin=*, nosep]
			\item In the actions drop-down, select \button{PickSurface}.
			\item Hold \key{Ctrl} and click \key{LMB} on the surface in the 3D view.
		\end{enumerate} \\
		\midrule
		\textbf{Expected Result} &
		\begin{itemize}[leftmargin=*, nosep]
			\item The surface's name turns green in the \tab{Surfaces} list.
			\item The surface's properties appear in the Properties panel.
		\end{itemize} \\
		\midrule
		\textbf{Result} & \checkbox~Pass \quad \checkbox~Fail \\
		\midrule
		\textbf{Notes} & \\[1.5cm]
		\bottomrule
	\end{tabularx}
	
	%----------------------------------------------------------------------------------------
	% SECTION 4: SURFACE PROPERTIES
	%----------------------------------------------------------------------------------------
	\newpage
	\section{Surface Properties and Controls}
	
	%--- tc:21 ---
	\testcase{21}{Surface Properties Panel}
	
	\begin{tabularx}{\textwidth}{lX}
		\toprule
		\textbf{Prerequisites} & A surface is selected (\ref{tc:19} passed). \\
		\midrule
		\textbf{Steps} &
		\begin{enumerate}[leftmargin=*, nosep]
			\item With the surface selected, view the Properties panel on the right.
			\item Change the surface name by clicking the name field, typing a new name, and pressing \key{Enter}.
			\item Change the \menu{Priority} value to \texttt{0.1} and press \key{Enter}.
			\item Change the \menu{FillMode} to \texttt{Wireframe}.
		\end{enumerate} \\
		\midrule
		\textbf{Expected Result} &
		\begin{itemize}[leftmargin=*, nosep]
			\item The new name appears in the \tab{Surfaces} list.
			\item The priority change is accepted (no error).
			\item The surface renders as a wireframe mesh in the 3D view.
		\end{itemize} \\
		\midrule
		\textbf{Result} & \checkbox~Pass \quad \checkbox~Fail \\
		\midrule
		\textbf{Notes} & \\[1.5cm]
		\bottomrule
	\end{tabularx}
	
	\vspace{1em}
	
	%--- tc:22 ---
	\testcase{22}{Surface Visibility Toggle}
	
	\begin{tabularx}{\textwidth}{lX}
		\toprule
		\textbf{Prerequisites} & A surface is loaded and visible in the 3D view. \\
		\midrule
		\textbf{Steps} &
		\begin{enumerate}[leftmargin=*, nosep]
			\item In the \tab{Surfaces} list, click \button{Toggle Visible} below the surface name.
			\item Observe the 3D view.
			\item Click \button{Toggle Visible} again.
			\item Observe the 3D view.
		\end{enumerate} \\
		\midrule
		\textbf{Expected Result} &
		\begin{itemize}[leftmargin=*, nosep]
			\item After the first click, the surface disappears from the 3D view.
			\item After the second click, the surface reappears.
		\end{itemize} \\
		\midrule
		\textbf{Result} & \checkbox~Pass \quad \checkbox~Fail \\
		\midrule
		\textbf{Notes} & \\[1.5cm]
		\bottomrule
	\end{tabularx}
	
	\vspace{1em}
	
	%--- tc:23 ---
	\testcase{23}{Surface Color Correction}
	
	\begin{tabularx}{\textwidth}{lX}
		\toprule
		\textbf{Prerequisites} & A surface is selected and visible (\ref{tc:19} passed). \\
		\midrule
		\textbf{Steps} &
		\begin{enumerate}[leftmargin=*, nosep]
			\item With the surface selected, locate the \menu{Color Correction} panel in the Properties section.
			\item Adjust the \menu{Contrast} slider.
			\item Adjust the \menu{Brightness} slider.
			\item Observe the 3D view after each change.
		\end{enumerate} \\
		\midrule
		\textbf{Expected Result} &
		\begin{itemize}[leftmargin=*, nosep]
			\item Adjusting the contrast slider visibly changes the contrast of the surface texture in the 3D view.
			\item Adjusting the brightness slider visibly lightens or darkens the surface texture.
		\end{itemize} \\
		\midrule
		\textbf{Result} & \checkbox~Pass \quad \checkbox~Fail \\
		\midrule
		\textbf{Notes} & \\[1.5cm]
		\bottomrule
	\end{tabularx}
	
	\vspace{1em}
	
	%--- tc:24 ---
	\testcase{24}{Surface Translation}
	
	\begin{tabularx}{\textwidth}{lX}
		\toprule
		\textbf{Prerequisites} & A surface is selected and a Coordinate System has been placed (\ref{tc:10} passed). \\
		\midrule
		\textbf{Steps} &
		\begin{enumerate}[leftmargin=*, nosep]
			\item With the surface selected, locate the \menu{Translation} panel in the Properties section.
			\item Enter a small offset value (e.g., \texttt{1.0}) in the \menu{North} axis field and press \key{Enter}.
			\item Observe the 3D view.
			\item Reset the value back to \texttt{0.0}.
		\end{enumerate} \\
		\midrule
		\textbf{Expected Result} &
		\begin{itemize}[leftmargin=*, nosep]
			\item The surface moves by the specified amount along the north axis.
			\item The 3D view updates immediately.
			\item Resetting to 0 returns the surface to its original position.
		\end{itemize} \\
		\midrule
		\textbf{Result} & \checkbox~Pass \quad \checkbox~Fail \\
		\midrule
		\textbf{Notes} & \\[1.5cm]
		\bottomrule
	\end{tabularx}
	
	\vspace{1em}
	
	%--- tc:25 ---
	\testcase{25}{Surface FillMode (Wireframe)}
	
	\begin{tabularx}{\textwidth}{lX}
		\toprule
		\textbf{Prerequisites} & A surface is selected. \\
		\midrule
		\textbf{Steps} &
		\begin{enumerate}[leftmargin=*, nosep]
			\item In the surface Properties panel, find the \menu{FillMode} drop-down.
			\item Set it to \texttt{Wireframe}.
			\item Observe the 3D view.
			\item Set it to \texttt{Point}.
			\item Observe the 3D view.
			\item Set it back to \texttt{Solid}.
		\end{enumerate} \\
		\midrule
		\textbf{Expected Result} &
		\begin{itemize}[leftmargin=*, nosep]
			\item \textbf{Wireframe}: The surface is rendered as a mesh of triangles (edges only).
			\item \textbf{Point}: The surface is rendered as a point cloud (vertices only).
			\item \textbf{Solid}: The surface renders normally with filled triangles and texture.
		\end{itemize} \\
		\midrule
		\textbf{Result} & \checkbox~Pass \quad \checkbox~Fail \\
		\midrule
		\textbf{Notes} & \\[1.5cm]
		\bottomrule
	\end{tabularx}
	
	%----------------------------------------------------------------------------------------
	% SECTION 5: ANNOTATION PROPERTIES
	%----------------------------------------------------------------------------------------
	\newpage
	\section{Annotation Properties and Measurements}
	
	%--- tc:26 ---
	\testcase{26}{Annotation Properties Panel}
	
	\begin{tabularx}{\textwidth}{lX}
		\toprule
		\textbf{Prerequisites} & A line or polyline annotation is selected (\ref{tc:17} passed). \\
		\midrule
		\textbf{Steps} &
		\begin{enumerate}[leftmargin=*, nosep]
			\item With an annotation selected, find the Properties panel.
			\item Change the \menu{Color} to a different color.
			\item Change the \menu{Thickness} to a larger value (e.g., \texttt{5}).
			\item Toggle the \menu{Visible} checkbox off, then on.
			\item Change the \menu{TextSize} value.
		\end{enumerate} \\
		\midrule
		\textbf{Expected Result} &
		\begin{itemize}[leftmargin=*, nosep]
			\item The annotation updates its color in the 3D view immediately.
			\item The line thickness visibly increases in the 3D view.
			\item Toggling visibility hides and then shows the annotation in the 3D view.
			\item Text size changes affect any label displayed near the annotation.
		\end{itemize} \\
		\midrule
		\textbf{Result} & \checkbox~Pass \quad \checkbox~Fail \\
		\midrule
		\textbf{Notes} & \\[1.5cm]
		\bottomrule
	\end{tabularx}
	
	\vspace{1em}
	
	%--- tc:27 ---
	\testcase{27}{Annotation Measurements}
	
	\begin{tabularx}{\textwidth}{lX}
		\toprule
		\textbf{Prerequisites} & A polyline annotation with multiple points is selected. \\
		\midrule
		\textbf{Steps} &
		\begin{enumerate}[leftmargin=*, nosep]
			\item With a polyline annotation selected, navigate to the \menu{Measurements} panel.
			\item Read the following values: \menu{Length}, \menu{WayLength}, \menu{Height}, \menu{HeightDelta}, \menu{Slope}, \menu{Bearing}.
		\end{enumerate} \\
		\midrule
		\textbf{Expected Result} &
		\begin{itemize}[leftmargin=*, nosep]
			\item All measurement fields are populated with numeric values.
			\item \menu{Length} shows the straight-line distance sum.
			\item \menu{WayLength} shows the projected distance along the surface.
			\item \menu{Height} shows the elevation difference between start and end.
			\item \menu{Slope} and \menu{Bearing} are displayed in appropriate units (degrees).
		\end{itemize} \\
		\midrule
		\textbf{Result} & \checkbox~Pass \quad \checkbox~Fail \\
		\midrule
		\textbf{Notes} & \\[1.5cm]
		\bottomrule
	\end{tabularx}
	
	\vspace{1em}
	
	%--- tc:28 ---
	\testcase{28}{Annotation Text Note}
	
	\begin{tabularx}{\textwidth}{lX}
		\toprule
		\textbf{Prerequisites} & An annotation is selected. \\
		\midrule
		\textbf{Steps} &
		\begin{enumerate}[leftmargin=*, nosep]
			\item In the annotation Properties panel, locate the \menu{Text} field.
			\item Type a short note (e.g., \texttt{Test Note}).
			\item Press \key{Enter}.
			\item Observe the 3D view.
		\end{enumerate} \\
		\midrule
		\textbf{Expected Result} &
		\begin{itemize}[leftmargin=*, nosep]
			\item The entered text appears as a label next to the annotation in the 3D view.
		\end{itemize} \\
		\midrule
		\textbf{Result} & \checkbox~Pass \quad \checkbox~Fail \\
		\midrule
		\textbf{Notes} & \\[1.5cm]
		\bottomrule
	\end{tabularx}
	
	%----------------------------------------------------------------------------------------
	% SECTION 6: SCALEBARS
	%----------------------------------------------------------------------------------------
	\newpage
	\section{Scalebars}
	
	%--- tc:29 ---
	\testcase{29}{Place Scalebar}
	
	\begin{tabularx}{\textwidth}{lX}
		\toprule
		\textbf{Prerequisites} & A surface is loaded. A coordinate system has been placed (\ref{tc:10} passed). \\
		\midrule
		\textbf{Steps} &
		\begin{enumerate}[leftmargin=*, nosep]
			\item In the actions drop-down, select \button{PlaceScaleBar}.
			\item Select an orientation mode (e.g., \texttt{Horizontal\_cam}).
			\item Choose a position (e.g., \texttt{Middle}).
			\item Enter a length value (e.g., \texttt{10}) and select a unit (e.g., \texttt{m}).
			\item Hold \key{Ctrl} and click \key{LMB} on the surface to place the scalebar.
		\end{enumerate} \\
		\midrule
		\textbf{Expected Result} &
		\begin{itemize}[leftmargin=*, nosep]
			\item A scalebar appears on the surface in the 3D view, aligned according to the selected orientation.
			\item The length label (e.g., \texttt{10 m}) is shown beside the scalebar.
			\item The scalebar is listed in the \tab{ScaleBars} panel.
		\end{itemize} \\
		\midrule
		\textbf{Result} & \checkbox~Pass \quad \checkbox~Fail \\
		\midrule
		\textbf{Notes} & \\[1.5cm]
		\bottomrule
	\end{tabularx}
	
	\vspace{1em}
	
	%--- tc:30 ---
	\testcase{30}{Scalebar Properties}
	
	\begin{tabularx}{\textwidth}{lX}
		\toprule
		\textbf{Prerequisites} & A scalebar exists (\ref{tc:29} passed) and is selected. \\
		\midrule
		\textbf{Steps} &
		\begin{enumerate}[leftmargin=*, nosep]
			\item Click the scalebar name in the \tab{ScaleBars} list to select it.
			\item Change the \menu{Thickness} value.
			\item Change the \menu{Subdivisions} value (e.g., to \texttt{5}).
			\item Toggle \menu{TextVisible} to hide the label, then show it again.
			\item Click the \button{Toggle} button below the scalebar name to hide/show it.
			\item Click the \button{remove} button below the scalebar name to delete it.
		\end{enumerate} \\
		\midrule
		\textbf{Expected Result} &
		\begin{itemize}[leftmargin=*, nosep]
			\item Thickness change updates the scalebar line width visually.
			\item Subdivisions change updates the number of black-white segments on the bar.
			\item Toggling text visibility hides/shows the length label.
			\item Clicking Toggle hides/shows the entire scalebar.
			\item Clicking remove deletes the scalebar from the list and the 3D view.
		\end{itemize} \\
		\midrule
		\textbf{Result} & \checkbox~Pass \quad \checkbox~Fail \\
		\midrule
		\textbf{Notes} & \\[1.5cm]
		\bottomrule
	\end{tabularx}
	
	%----------------------------------------------------------------------------------------
	% SECTION 7: BOOKMARKS
	%----------------------------------------------------------------------------------------
	\newpage
	\section{Bookmarks}
	
	%--- tc:31 ---
	\testcase{31}{Add Bookmark}
	
	\begin{tabularx}{\textwidth}{lX}
		\toprule
		\textbf{Prerequisites} & A surface is loaded. Navigate to a desired viewpoint in the 3D view. \\
		\midrule
		\textbf{Steps} &
		\begin{enumerate}[leftmargin=*, nosep]
			\item Navigate to the \tab{Bookmarks} page in the right panel.
			\item Click the \button{+} button at the top of the bookmarks page.
		\end{enumerate} \\
		\midrule
		\textbf{Expected Result} &
		\begin{itemize}[leftmargin=*, nosep]
			\item A new bookmark entry appears in the bookmarks list.
			\item The bookmark captures the current camera position.
		\end{itemize} \\
		\midrule
		\textbf{Result} & \checkbox~Pass \quad \checkbox~Fail \\
		\midrule
		\textbf{Notes} & \\[1.5cm]
		\bottomrule
	\end{tabularx}
	
	\vspace{1em}
	
	%--- tc:32 ---
	\testcase{32}{Bookmark FlyTo}
	
	\begin{tabularx}{\textwidth}{lX}
		\toprule
		\textbf{Prerequisites} & At least one bookmark exists (\ref{tc:31} passed). Move the camera to a different viewpoint. \\
		\midrule
		\textbf{Steps} &
		\begin{enumerate}[leftmargin=*, nosep]
			\item In the \tab{Bookmarks} list, click the ``house'' icon beside the bookmark's name.
		\end{enumerate} \\
		\midrule
		\textbf{Expected Result} &
		\begin{itemize}[leftmargin=*, nosep]
			\item The camera animates (flies) back to the saved viewpoint.
			\item The 3D view shows the scene from the position saved in the bookmark.
		\end{itemize} \\
		\midrule
		\textbf{Result} & \checkbox~Pass \quad \checkbox~Fail \\
		\midrule
		\textbf{Notes} & \\[1.5cm]
		\bottomrule
	\end{tabularx}
	
	%----------------------------------------------------------------------------------------
	% SECTION 8: SEQUENCED BOOKMARKS
	%----------------------------------------------------------------------------------------
	\newpage
	\section{Sequenced Bookmarks}
	
	%--- tc:33 ---
	\testcase{33}{Sequenced Bookmarks -- Add and Animate}
	
	\begin{tabularx}{\textwidth}{lX}
		\toprule
		\textbf{Prerequisites} & A surface is loaded. Navigate to at least two different viewpoints to capture. \\
		\midrule
		\textbf{Steps} &
		\begin{enumerate}[leftmargin=*, nosep]
			\item Navigate to the \tab{Sequenced Bookmarks} page.
			\item Navigate the camera to a first position, then click the \button{+} button (A in the UI) to add a bookmark.
			\item Navigate the camera to a second position, then click \button{+} again.
			\item In the bookmark list (B), set a \menu{Duration} for the second bookmark (e.g., \texttt{3} seconds).
			\item In the Animation section (D), click the \textbf{Play} button ($\blacktriangleright$).
		\end{enumerate} \\
		\midrule
		\textbf{Expected Result} &
		\begin{itemize}[leftmargin=*, nosep]
			\item Two bookmarks appear in the sequenced bookmarks list.
			\item Clicking Play starts a camera animation that travels from the first to the second bookmark.
			\item The animation duration corresponds to the set value.
			\item The Pause and Stop buttons become active during animation.
		\end{itemize} \\
		\midrule
		\textbf{Result} & \checkbox~Pass \quad \checkbox~Fail \\
		\midrule
		\textbf{Notes} & \\[1.5cm]
		\bottomrule
	\end{tabularx}
	
	\vspace{1em}
	
	%--- tc:34 ---
	\testcase{34}{Sequenced Bookmarks -- Record and Generate Images}
	
	\begin{tabularx}{\textwidth}{lX}
		\toprule
		\textbf{Prerequisites} & At least two sequenced bookmarks exist (\ref{tc:33} passed). An output path is set in the Snapshots section (E). \\
		\midrule
		\textbf{Steps} &
		\begin{enumerate}[leftmargin=*, nosep]
			\item In the \menu{Snapshots} section (E), ensure the \menu{Output Path} is valid.
			\item Set an \menu{Image Resolution} (e.g., \texttt{1920x1080}).
			\item Click the red \textbf{Record} button in the Snapshots section.
			\item Click the \textbf{Play} button in the Animation section (D) to animate the camera.
			\item Wait for the animation to complete.
			\item Click the red \textbf{Stop} button in the Snapshots section.
			\item Click the \textbf{Generate Images} button (camera icon).
		\end{enumerate} \\
		\midrule
		\textbf{Expected Result} &
		\begin{itemize}[leftmargin=*, nosep]
			\item The Record button changes to a Stop button when recording starts.
			\item After stopping, a JSON file is saved in the output folder.
			\item After clicking Generate Images, PRo3D renders images in the background.
			\item Images (named \texttt{img\_000.png} etc.) appear in the output folder.
		\end{itemize} \\
		\midrule
		\textbf{Result} & \checkbox~Pass \quad \checkbox~Fail \\
		\midrule
		\textbf{Notes} & \\[1.5cm]
		\bottomrule
	\end{tabularx}
	
	%----------------------------------------------------------------------------------------
	% SECTION 9: VIEWER CONFIGURATION
	%----------------------------------------------------------------------------------------
	\newpage
	\section{Viewer Configuration}
	
	%--- tc:35 ---
	\testcase{35}{Viewer Config Settings}
	
	\begin{tabularx}{\textwidth}{lX}
		\toprule
		\textbf{Prerequisites} & A surface is loaded. \\
		\midrule
		\textbf{Steps} &
		\begin{enumerate}[leftmargin=*, nosep]
			\item Navigate to the \tab{Config} page in the right panel.
			\item Adjust the \menu{Navigation Sensitivity} slider.
			\item Change the \menu{Arrow Length} value.
			\item Enable \menu{Lod Colors} (if available) and observe.
			\item Press \key{Page Up} to increase navigation sensitivity.
			\item Press \key{Page Down} to decrease navigation sensitivity.
		\end{enumerate} \\
		\midrule
		\textbf{Expected Result} &
		\begin{itemize}[leftmargin=*, nosep]
			\item Adjusting the Navigation Sensitivity slider changes how fast the camera responds to input.
			\item \key{Page Up}/\key{Page Down} also adjusts sensitivity (the value updates in the Config panel).
			\item Arrow Length change affects the size of direction arrows visible in the scene.
			\item Lod Colors colors the geometry by level-of-detail (different shades visible on the surface).
		\end{itemize} \\
		\midrule
		\textbf{Result} & \checkbox~Pass \quad \checkbox~Fail \\
		\midrule
		\textbf{Notes} & \\[1.5cm]
		\bottomrule
	\end{tabularx}
	
	\vspace{1em}
	
	%--- tc:36 ---
	\testcase{36}{Screenshots}
	
	\begin{tabularx}{\textwidth}{lX}
		\toprule
		\textbf{Prerequisites} & A surface is loaded and visible. \\
		\midrule
		\textbf{Steps} &
		\begin{enumerate}[leftmargin=*, nosep]
			\item Navigate to the \tab{Config} page.
			\item Locate the \menu{Screenshots} section.
			\item Set a width and height (e.g., \texttt{1920} and \texttt{1080}).
			\item Select output format \texttt{PNG}.
			\item Click the \textbf{camera icon} button to capture a screenshot.
			\item Click the \textbf{folder icon} button to open the output folder.
		\end{enumerate} \\
		\midrule
		\textbf{Expected Result} &
		\begin{itemize}[leftmargin=*, nosep]
			\item Clicking the camera icon creates a screenshot file in the output directory.
			\item Clicking the folder icon opens the operating system's file explorer showing the output folder.
			\item The screenshot file (e.g., \texttt{img0.png}) is present in the folder and matches the current 3D view.
		\end{itemize} \\
		\midrule
		\textbf{Result} & \checkbox~Pass \quad \checkbox~Fail \\
		\midrule
		\textbf{Notes} & \\[1.5cm]
		\bottomrule
	\end{tabularx}
	
	%----------------------------------------------------------------------------------------
	% SECTION 10: GROUPING
	%----------------------------------------------------------------------------------------
	\newpage
	\section{Grouping}
	
	%--- tc:37 ---
	\testcase{37}{Grouping -- Create Group}
	
	\begin{tabularx}{\textwidth}{lX}
		\toprule
		\textbf{Prerequisites} & At least two annotations exist. \\
		\midrule
		\textbf{Steps} &
		\begin{enumerate}[leftmargin=*, nosep]
			\item Navigate to the \tab{Annotations} panel.
			\item In the listing, right-click or use the group context menu on the ``root'' group.
			\item Select \button{Add Group} from the context menu.
			\item Verify a new empty group appears below root.
			\item Right-click the new group and select \button{Set Active}.
		\end{enumerate} \\
		\midrule
		\textbf{Expected Result} &
		\begin{itemize}[leftmargin=*, nosep]
			\item A new subgroup entry appears in the annotations hierarchy.
			\item Setting it as active means any new annotation will be placed in this group.
		\end{itemize} \\
		\midrule
		\textbf{Result} & \checkbox~Pass \quad \checkbox~Fail \\
		\midrule
		\textbf{Notes} & \\[1.5cm]
		\bottomrule
	\end{tabularx}
	
	\vspace{1em}
	
	%--- tc:38 ---
	\testcase{38}{Grouping -- Move Leafs}
	
	\begin{tabularx}{\textwidth}{lX}
		\toprule
		\textbf{Prerequisites} & At least two annotations exist, and a group has been created (\ref{tc:37} passed). \\
		\midrule
		\textbf{Steps} &
		\begin{enumerate}[leftmargin=*, nosep]
			\item In the \tab{Annotations} list, click the square icon in front of one or more annotations. The square turns green.
			\item Set the target group as active (right-click $\rightarrow$ \button{Set Active}).
			\item In the leaf actions panel, click \button{Selection: Move}.
		\end{enumerate} \\
		\midrule
		\textbf{Expected Result} &
		\begin{itemize}[leftmargin=*, nosep]
			\item The selected annotations (green squares) move to the active group in the hierarchy.
			\item They no longer appear under the root group, but appear under the target group.
		\end{itemize} \\
		\midrule
		\textbf{Result} & \checkbox~Pass \quad \checkbox~Fail \\
		\midrule
		\textbf{Notes} & \\[1.5cm]
		\bottomrule
	\end{tabularx}
	
	%----------------------------------------------------------------------------------------
	% SECTION 11: VIEW PLANNER
	%----------------------------------------------------------------------------------------
	\newpage
	\section{View Planner}
	
	%--- tc:39 ---
	\testcase{39}{View Planner -- Place Rover}
	
	\begin{tabularx}{\textwidth}{lX}
		\toprule
		\textbf{Prerequisites} & A surface is loaded. A \texttt{rover.xml} file is present in the \texttt{Release\textbackslash InstrumentStuff} folder. \\
		\midrule
		\textbf{Steps} &
		\begin{enumerate}[leftmargin=*, nosep]
			\item In the actions drop-down, select \button{PlaceRover}.
			\item From the rover model selector, choose a rover model.
			\item Hold \key{Ctrl} and click \key{LMB} on a first point on the surface (rover position -- shown in green).
			\item Hold \key{Ctrl} and click \key{LMB} on a second point on the surface (viewing direction -- shown in yellow).
		\end{enumerate} \\
		\midrule
		\textbf{Expected Result} &
		\begin{itemize}[leftmargin=*, nosep]
			\item A rover model appears on the surface at the first picked point.
			\item The rover faces the direction indicated by the second point.
			\item The View Plan entry appears in the \tab{ViewPlanner} list.
		\end{itemize} \\
		\midrule
		\textbf{Result} & \checkbox~Pass \quad \checkbox~Fail \\
		\midrule
		\textbf{Notes} & \\[1.5cm]
		\bottomrule
	\end{tabularx}
	
	\vspace{1em}
	
	%--- tc:40 ---
	\testcase{40}{View Planner -- Instrument View}
	
	\begin{tabularx}{\textwidth}{lX}
		\toprule
		\textbf{Prerequisites} & A rover has been placed (\ref{tc:39} passed). \\
		\midrule
		\textbf{Steps} &
		\begin{enumerate}[leftmargin=*, nosep]
			\item In the \tab{ViewPlanner} list, click the square icon next to a View Plan to select it.
			\item From the \menu{Instrument} drop-down, select a camera instrument.
			\item Adjust \menu{Sensor (px)} and \menu{Focal (mm)} parameters.
			\item Adjust \menu{Pan} and \menu{Tilt} angles.
			\item Enable \menu{show footprint}.
			\item Click \menu{export footprint}.
		\end{enumerate} \\
		\midrule
		\textbf{Expected Result} &
		\begin{itemize}[leftmargin=*, nosep]
			\item The camera footprint (light gray projection) appears on the surface in the 3D view.
			\item An Instrument View panel appears showing the view from the selected instrument.
			\item Adjusting pan/tilt updates the footprint and instrument view.
			\item Exporting footprint saves a screenshot from the main view, an instrument view screenshot, and a \texttt{.svx} metadata file.
		\end{itemize} \\
		\midrule
		\textbf{Result} & \checkbox~Pass \quad \checkbox~Fail \\
		\midrule
		\textbf{Notes} & \\[1.5cm]
		\bottomrule
	\end{tabularx}
	
	%----------------------------------------------------------------------------------------
	% SECTION 12: GIS VIEW
	%----------------------------------------------------------------------------------------
	\newpage
	\section{GIS View}
	
	%--- tc:41 ---
	\testcase{41}{GIS View -- Load SPICE Kernel}
	
	\begin{tabularx}{\textwidth}{lX}
		\toprule
		\textbf{Prerequisites} & GIS View is available (application started with GIS mode enabled or via Change Mode). A SPICE kernel file (\texttt{.tm}) is available on disk. \\
		\midrule
		\textbf{Steps} &
		\begin{enumerate}[leftmargin=*, nosep]
			\item Navigate to the \tab{GIS} tab in the right panel.
			\item Locate the kernel loading section and click the button to browse for or enter the kernel path.
			\item Select or enter the path to the \texttt{.tm} kernel file.
			\item Confirm the selection.
		\end{enumerate} \\
		\midrule
		\textbf{Expected Result} &
		\begin{itemize}[leftmargin=*, nosep]
			\item The kernel path appears in the GIS View settings pane.
			\item A green check mark appears below the path, indicating successful loading.
			\item If loading fails, a red exclamation mark is shown (check PRo3D console for details).
		\end{itemize} \\
		\midrule
		\textbf{Result} & \checkbox~Pass \quad \checkbox~Fail \\
		\midrule
		\textbf{Notes} & \\[1.5cm]
		\bottomrule
	\end{tabularx}
	
	\vspace{1em}
	
	%--- tc:42 ---
	\testcase{42}{GIS View -- Observation Settings}
	
	\begin{tabularx}{\textwidth}{lX}
		\toprule
		\textbf{Prerequisites} & A SPICE kernel is loaded (\ref{tc:41} passed). \\
		\midrule
		\textbf{Steps} &
		\begin{enumerate}[leftmargin=*, nosep]
			\item In the \menu{Current Observation Settings} section, set an \menu{Observed Body} (e.g., \texttt{MARS}).
			\item Set a \menu{Camera Source Body} (e.g., \texttt{PHOBOS}).
			\item Set a valid \menu{Time} (within the kernel's coverage).
			\item Set a \menu{Reference Frame}.
			\item Click the \button{Reset} button.
		\end{enumerate} \\
		\midrule
		\textbf{Expected Result} &
		\begin{itemize}[leftmargin=*, nosep]
			\item The camera moves to the location of the Camera Source Body and points toward the Observed Body.
			\item The 3D view updates to show the configured observation scenario.
			\item Entities with proxy geometry enabled are visible in the scene.
		\end{itemize} \\
		\midrule
		\textbf{Result} & \checkbox~Pass \quad \checkbox~Fail \\
		\midrule
		\textbf{Notes} & \\[1.5cm]
		\bottomrule
	\end{tabularx}
	
	%----------------------------------------------------------------------------------------
	% SECTION 13: CONTOUR LINES AND MULTITEXTURING
	%----------------------------------------------------------------------------------------
	\newpage
	\section{Contour Lines and Multitexturing}
	
	%--- tc:43 ---
	\testcase{43}{Contour Lines}
	
	\begin{tabularx}{\textwidth}{lX}
		\toprule
		\textbf{Prerequisites} & A multilayer OPC surface (with an \texttt{.opcx} file and an elevation layer) is loaded and selected. \\
		\midrule
		\textbf{Steps} &
		\begin{enumerate}[leftmargin=*, nosep]
			\item With the surface selected, in the surface Properties panel, navigate to the \menu{Textures} selector and choose the elevation layer as the secondary texture.
			\item Find the \menu{Contour} section in the surface properties.
			\item Enable contour lines by checking the enable checkbox.
			\item Set a \menu{Distance} value (line interval).
			\item Set \menu{Line Width} and \menu{Line Border} values.
		\end{enumerate} \\
		\midrule
		\textbf{Expected Result} &
		\begin{itemize}[leftmargin=*, nosep]
			\item Contour lines appear on the surface in the 3D view, spaced by the specified distance.
			\item Adjusting the interval changes the spacing between contour lines visibly.
		\end{itemize} \\
		\midrule
		\textbf{Result} & \checkbox~Pass \quad \checkbox~Fail \\
		\midrule
		\textbf{Notes} & \\[1.5cm]
		\bottomrule
	\end{tabularx}
	
	\vspace{1em}
	
	%--- tc:44 ---
	\testcase{44}{Multitexturing}
	
	\begin{tabularx}{\textwidth}{lX}
		\toprule
		\textbf{Prerequisites} & A surface with an \texttt{.opcx} file providing multiple texture and/or scalar layers is loaded and selected. \\
		\midrule
		\textbf{Steps} &
		\begin{enumerate}[leftmargin=*, nosep]
			\item With the surface selected, find the \menu{Scalars} and \menu{Textures} combo boxes in the Properties panel.
			\item Select a different layer from the \menu{Scalars} drop-down (e.g., an elevation or accuracy layer).
			\item Locate the \menu{Multitexturing} controls in the surface properties.
			\item Adjust the opacity or blending slider.
		\end{enumerate} \\
		\midrule
		\textbf{Expected Result} &
		\begin{itemize}[leftmargin=*, nosep]
			\item The selected scalar/attribute layer is applied as a false-color map on the surface.
			\item A color legend appears on the left side of the 3D view.
			\item Adjusting the blending slider changes how much the additional layer overlays the base texture.
		\end{itemize} \\
		\midrule
		\textbf{Result} & \checkbox~Pass \quad \checkbox~Fail \\
		\midrule
		\textbf{Notes} & \\[1.5cm]
		\bottomrule
	\end{tabularx}
	
	\vspace{1em}
	
	%--- tc:57 ---
	\testcase{57}{Cross Sections}
	
	\begin{tabularx}{\textwidth}{lX}
		\toprule
		\textbf{Prerequisites} & A surface is loaded and visible. A curtain image (PNG) prepared from an exported profile is available on disk (see the \menu{Cross Sections} section of the Short User Manual). \\
		\midrule
		\textbf{Steps} &
		\begin{enumerate}[leftmargin=*, nosep]
			\item In \button{DrawAnnotation} mode, choose \button{Polyline} with \button{Sky} projection and draw a polyline along the area to be cut. Use an adequate sampling distance to avoid excessive segment counts.
			\item Select the polyline annotation and open its Properties panel.
			\item Use \button{Export Profile} to export the profile as a CSV file.
			\item Move the camera far away from the cross section so that everything between the annotation and the camera will be clipped.
			\item In the annotation properties, trigger \button{Create Cross Section}.
			\item Open the \menu{Curtain} settings, set the curtain height and minimum altitude to match the values used when generating the curtain image, and select the prepared curtain image.
		\end{enumerate} \\
		\midrule
		\textbf{Expected Result} &
		\begin{itemize}[leftmargin=*, nosep]
			\item The exported CSV contains \texttt{distance,elevation} columns with one row per sampled point.
			\item After creating the cross section, the surface data in front of the cutting plane is clipped away, leaving the cut visible.
			\item The curtain image is mapped onto the cutting plane and aligned with the profile.
			\item The cross section and curtain can be inspected in the 3D view without artifacts.
		\end{itemize} \\
		\midrule
		\textbf{Result} & \checkbox~Pass \quad \checkbox~Fail \\
		\midrule
		\textbf{Notes} & \\[1.5cm]
		\bottomrule
	\end{tabularx}
	
	%----------------------------------------------------------------------------------------
	% SECTION 14: SURFACE COMPARISON
	%----------------------------------------------------------------------------------------
	\newpage
	\section{Surface Comparison}
	
	%--- tc:45 ---
	\testcase{45}{Surface Comparison -- Activate Mode}
	
	\begin{tabularx}{\textwidth}{lX}
		\toprule
		\textbf{Prerequisites} & Two surfaces with different names are loaded. \\
		\midrule
		\textbf{Steps} &
		\begin{enumerate}[leftmargin=*, nosep]
			\item Click the main menu icon.
			\item Navigate to \menu{Change Mode $\rightarrow$ Surface Comparison}.
			\item In the Surface Comparison tab, select one surface from the \menu{Surface1} drop-down.
			\item Select the second surface from the \menu{Surface2} drop-down.
			\item Press \key{t} on the keyboard.
		\end{enumerate} \\
		\midrule
		\textbf{Expected Result} &
		\begin{itemize}[leftmargin=*, nosep]
			\item The Surface Comparison interface appears with Surface1 and Surface2 selectors.
			\item Both surfaces can be selected independently.
			\item Pressing \key{t} toggles visibility between the two surfaces (only one is visible at a time).
		\end{itemize} \\
		\midrule
		\textbf{Result} & \checkbox~Pass \quad \checkbox~Fail \\
		\midrule
		\textbf{Notes} & \\[1.5cm]
		\bottomrule
	\end{tabularx}
	
	\vspace{1em}
	
	%--- tc:46 ---
	\testcase{46}{Surface Comparison -- Select Area}
	
	\begin{tabularx}{\textwidth}{lX}
		\toprule
		\textbf{Prerequisites} & Surface Comparison mode is active and two surfaces are selected (\ref{tc:45} passed). \\
		\midrule
		\textbf{Steps} &
		\begin{enumerate}[leftmargin=*, nosep]
			\item In the menu bar above the 3D view, choose \button{Select Area} from the drop-down.
			\item Hold \key{Ctrl} and click \key{LMB} on the surface to select a location.
			\item Observe the sphere indicator.
			\item Press \key{+} to increase the sphere size; press \key{-} to decrease it.
			\item Click the calculator icon button to update measurements.
		\end{enumerate} \\
		\midrule
		\textbf{Expected Result} &
		\begin{itemize}[leftmargin=*, nosep]
			\item A translucent sphere appears at the clicked location.
			\item The \key{+}/\key{-} keys change the sphere radius.
			\item After clicking the calculator icon, colored spheres appear within the area, color-coded by distance between the two surfaces (red = large difference, green = small difference).
			\item A color legend is shown on the left side.
			\item Statistics appear in the panel on the right for the selected area.
		\end{itemize} \\
		\midrule
		\textbf{Result} & \checkbox~Pass \quad \checkbox~Fail \\
		\midrule
		\textbf{Notes} & \\[1.5cm]
		\bottomrule
	\end{tabularx}
	
	\vspace{1em}
	
	%--- tc:47 ---
	\testcase{47}{Surface Comparison -- Surface Measurements}
	
	\begin{tabularx}{\textwidth}{lX}
		\toprule
		\textbf{Prerequisites} & Two surfaces are selected in the Comparison tab (\ref{tc:45} passed). \\
		\midrule
		\textbf{Steps} &
		\begin{enumerate}[leftmargin=*, nosep]
			\item In the Surface Comparison panel, navigate to the \menu{Surface Measurements} section.
			\item From the \menu{Origin} drop-down, select a center calculation method.
			\item Click the calculator icon button to update measurements.
		\end{enumerate} \\
		\midrule
		\textbf{Expected Result} &
		\begin{itemize}[leftmargin=*, nosep]
			\item Size measurements along the coordinate axes are listed for each surface.
			\item Differences between the two surfaces are shown below the individual measurements.
		\end{itemize} \\
		\midrule
		\textbf{Result} & \checkbox~Pass \quad \checkbox~Fail \\
		\midrule
		\textbf{Notes} & \\[1.5cm]
		\bottomrule
	\end{tabularx}
	
	\vspace{1em}
	
	%--- tc:48 ---
	\testcase{48}{Surface Comparison -- Length Measurements via Bookmarks}
	
	\begin{tabularx}{\textwidth}{lX}
		\toprule
		\textbf{Prerequisites} & Two surfaces are selected in Surface Comparison mode. A bookmark is available. \\
		\midrule
		\textbf{Steps} &
		\begin{enumerate}[leftmargin=*, nosep]
			\item Navigate to the \tab{Bookmarks} page and select a bookmark (it should be highlighted in green).
			\item Switch to \button{DrawAnnotation} mode (\key{F2}).
			\item In the projection selector, choose \button{Bookmark} projection mode.
			\item Draw one annotation on Surface1.
			\item Press \key{t} to switch to Surface2.
			\item Draw another annotation on Surface2 (same bookmark selected).
			\item Return to the Surface Comparison window and click the calculator icon.
		\end{enumerate} \\
		\midrule
		\textbf{Expected Result} &
		\begin{itemize}[leftmargin=*, nosep]
			\item The Bookmark projection mode is available in the DrawAnnotation panel.
			\item Both annotations use the same bookmark as the projection origin.
			\item The \menu{Annotation Length Comparison} section shows the length measurements for both annotations and their difference.
		\end{itemize} \\
		\midrule
		\textbf{Result} & \checkbox~Pass \quad \checkbox~Fail \\
		\midrule
		\textbf{Notes} & \\[1.5cm]
		\bottomrule
	\end{tabularx}
	
	%----------------------------------------------------------------------------------------
	% SECTION 15: MINERVA
	%----------------------------------------------------------------------------------------
	\newpage
	\section{Minerva}
	
	%--- tc:49 ---
	\testcase{49}{Minerva -- Load Data and Product Listing}
	
	\begin{tabularx}{\textwidth}{lX}
		\toprule
		\textbf{Prerequisites} & A \texttt{dump.csv} file is present in the \texttt{Release\textbackslash netcoreapp2.0\textbackslash MinervaData} folder. PRo3D is started. \\
		\midrule
		\textbf{Steps} &
		\begin{enumerate}[leftmargin=*, nosep]
			\item Start PRo3D.
			\item Navigate to the \tab{Minerva} tab in the right panel.
			\item Observe the product listing.
		\end{enumerate} \\
		\midrule
		\textbf{Expected Result} &
		\begin{itemize}[leftmargin=*, nosep]
			\item The Minerva tab is visible.
			\item Products are listed, grouped by instrument.
			\item Up to 20 products per instrument are shown in the list.
			\item Products appear as dots/markers in the 3D view after a surface is loaded.
		\end{itemize} \\
		\midrule
		\textbf{Result} & \checkbox~Pass \quad \checkbox~Fail \\
		\midrule
		\textbf{Notes} & \\[1.5cm]
		\bottomrule
	\end{tabularx}
	
	\vspace{1em}
	
	%--- tc:50 ---
	\testcase{50}{Minerva -- Query App Filter}
	
	\begin{tabularx}{\textwidth}{lX}
		\toprule
		\textbf{Prerequisites} & Minerva data is loaded (\ref{tc:49} passed). \\
		\midrule
		\textbf{Steps} &
		\begin{enumerate}[leftmargin=*, nosep]
			\item Click the \button{Query App} button (B in the Minerva tab) to open the Query App.
			\item Set a \menu{min sol} and \menu{max sol} range.
			\item Uncheck some instruments.
			\item Click \button{ApplyFilter}.
			\item Observe the product listing and 3D view.
			\item Click \button{ClearFilter}.
		\end{enumerate} \\
		\midrule
		\textbf{Expected Result} &
		\begin{itemize}[leftmargin=*, nosep]
			\item After applying the filter, only products within the sol range and checked instruments are shown.
			\item The 3D view updates to display only filtered products.
			\item Clicking ClearFilter restores all products.
		\end{itemize} \\
		\midrule
		\textbf{Result} & \checkbox~Pass \quad \checkbox~Fail \\
		\midrule
		\textbf{Notes} & \\[1.5cm]
		\bottomrule
	\end{tabularx}
	
	\vspace{1em}
	
	%--- tc:51 ---
	\testcase{51}{Minerva -- Product Selection}
	
	\begin{tabularx}{\textwidth}{lX}
		\toprule
		\textbf{Prerequisites} & Minerva data is loaded. Surface is loaded. \\
		\midrule
		\textbf{Steps} &
		\begin{enumerate}[leftmargin=*, nosep]
			\item In the product listing, click on a product name.
			\item Observe the Properties panel and the color change in the list.
			\item Click the \button{FlyTo} button below the product name.
			\item In the actions drop-down, set \button{PickMinervaProduct}.
			\item Hold \key{Ctrl} and click \key{LMB} on a product dot in the 3D view.
			\item Hold \key{Shift}+\key{LMB} and drag a rectangle over multiple products in the 3D view.
		\end{enumerate} \\
		\midrule
		\textbf{Expected Result} &
		\begin{itemize}[leftmargin=*, nosep]
			\item Clicking a product name highlights it in green in the list and shows its properties.
			\item FlyTo animates the camera to the product location.
			\item PickMinervaProduct allows selecting a product by clicking it in the 3D view.
			\item Shift-dragging selects multiple products within the drawn rectangle; all are highlighted.
		\end{itemize} \\
		\midrule
		\textbf{Result} & \checkbox~Pass \quad \checkbox~Fail \\
		\midrule
		\textbf{Notes} & \\[1.5cm]
		\bottomrule
	\end{tabularx}
	
	\vspace{1em}
	
	%--- tc:52 ---
	\testcase{52}{Minerva -- Pick Linking}
	
	\begin{tabularx}{\textwidth}{lX}
		\toprule
		\textbf{Prerequisites} & Minerva data is loaded. A surface is loaded. \\
		\midrule
		\textbf{Steps} &
		\begin{enumerate}[leftmargin=*, nosep]
			\item In the actions drop-down, select \button{PickLinking}.
			\item Hold \key{Ctrl} and click \key{LMB} to pick a point on the surface.
			\item Observe the 3D view.
		\end{enumerate} \\
		\midrule
		\textbf{Expected Result} &
		\begin{itemize}[leftmargin=*, nosep]
			\item Camera frustums are displayed for all products whose field of view captures the selected point.
			\item The frustums are drawn as wire cones/pyramids from each product's camera location toward the surface.
		\end{itemize} \\
		\midrule
		\textbf{Result} & \checkbox~Pass \quad \checkbox~Fail \\
		\midrule
		\textbf{Notes} & \\[1.5cm]
		\bottomrule
	\end{tabularx}
	
	%----------------------------------------------------------------------------------------
	% SECTION 16: COMMAND LINE INTERFACE
	%----------------------------------------------------------------------------------------
	\newpage
	\section{Command Line Interface}
	
	%--- tc:53 ---
	\testcase{53}{CLI -- Snapshot Rendering}
	
	\begin{tabularx}{\textwidth}{lX}
		\toprule
		\textbf{Prerequisites} & PRo3D is installed. An OPC surface folder is available. A valid snapshot JSON file (\texttt{snapshots.json}) is prepared according to the manual's format. \\
		\midrule
		\textbf{Steps} &
		\begin{enumerate}[leftmargin=*, nosep]
			\item Open a terminal (Command Prompt or PowerShell) and navigate to the directory containing \texttt{PRo3D.Snapshots.exe}.
			\item Run: \texttt{PRo3D.Snapshots.exe --help} and read the output.
			\item Run the basic snapshot command: \texttt{PRo3D.Snapshots.exe --opcs MyOPCs\textbackslash surface\textbackslash  --asnap snapshots.json --out MyImages\textbackslash }
			\item Wait for the rendering to complete.
			\item Check the output folder.
		\end{enumerate} \\
		\midrule
		\textbf{Expected Result} &
		\begin{itemize}[leftmargin=*, nosep]
			\item The \texttt{--help} flag prints a list of all available arguments.
			\item The rendering process starts and the viewer renders each snapshot.
			\item Rendered image files appear in the specified output folder (\texttt{MyImages\textbackslash}).
			\item The number of output images matches the number of snapshot entries in the JSON file.
		\end{itemize} \\
		\midrule
		\textbf{Result} & \checkbox~Pass \quad \checkbox~Fail \\
		\midrule
		\textbf{Notes} & \\[1.5cm]
		\bottomrule
	\end{tabularx}
	
	%----------------------------------------------------------------------------------------
	% SECTION 17: RIMFAX
	%----------------------------------------------------------------------------------------
	\newpage
	\section{RIMFAX}
	
	%--- tc:54 ---
	\testcase{54}{RIMFAX -- Load Traverse}
	
	\begin{tabularx}{\textwidth}{lX}
		\toprule
		\textbf{Prerequisites} & RIMFAX data (including a \texttt{RIMFAX\_traverse.json} file and the expected folder structure) is available. A surface is loaded. \\
		\midrule
		\textbf{Steps} &
		\begin{enumerate}[leftmargin=*, nosep]
			\item Click the main menu icon.
			\item Navigate to the \menu{Extra} panel.
			\item Click the button to load a RIMFAX traverse file and select the \texttt{RIMFAX\_traverse.json} file.
		\end{enumerate} \\
		\midrule
		\textbf{Expected Result} &
		\begin{itemize}[leftmargin=*, nosep]
			\item The RIMFAX traverse is loaded and displayed in the 3D view.
			\item The traverse path appears aligned with the rover path on the surface.
		\end{itemize} \\
		\midrule
		\textbf{Result} & \checkbox~Pass \quad \checkbox~Fail \\
		\midrule
		\textbf{Notes} & \\[1.5cm]
		\bottomrule
	\end{tabularx}
	
	\vspace{1em}
	
	%--- tc:55 ---
	\testcase{55}{RIMFAX -- Load Surfaces}
	
	\begin{tabularx}{\textwidth}{lX}
		\toprule
		\textbf{Prerequisites} & A RIMFAX traverse is loaded (\ref{tc:54} passed). The full RIMFAX folder structure (with OBJ/PNG files) is available. \\
		\midrule
		\textbf{Steps} &
		\begin{enumerate}[leftmargin=*, nosep]
			\item In the RIMFAX attribute window, click the \button{Import Surfaces} button.
			\item Navigate to the root RIMFAX data folder and select it.
			\item Observe the 3D view and the Sol panel.
		\end{enumerate} \\
		\midrule
		\textbf{Expected Result} &
		\begin{itemize}[leftmargin=*, nosep]
			\item RIMFAX surface images appear in the 3D view alongside the traverse.
			\item Sol numbers are listed in the Sol panel.
			\item The RIMFAX mode can be changed per sol in the Sol panel.
		\end{itemize} \\
		\midrule
		\textbf{Result} & \checkbox~Pass \quad \checkbox~Fail \\
		\midrule
		\textbf{Notes} & \\[1.5cm]
		\bottomrule
	\end{tabularx}
	
	%----------------------------------------------------------------------------------------
	% SECTION 18: KEYBOARD SHORTCUTS
	%----------------------------------------------------------------------------------------
	\newpage
	\section{Keyboard Shortcuts}
	
	%--- tc:56 ---
	\testcase{56}{Keyboard Shortcuts}
	
	\begin{tabularx}{\textwidth}{lX}
		\toprule
		\textbf{Prerequisites} & A surface is loaded. At least two surfaces are available for the surface comparison shortcut. \\
		\midrule
		\textbf{Steps} &
		\begin{enumerate}[leftmargin=*, nosep]
			\item Press \key{f} -- toggle linear texture magnification filtering. Observe the surface texture.
			\item Press \key{p} -- show/hide exploration point. Observe the pink explore center dot.
			\item Press \key{t} -- toggle visibility between two surfaces selected in the Comparison window.
			\item Press \key{Page Up} -- raise navigation sensitivity. Observe faster movement.
			\item Press \key{Page Down} -- lower navigation sensitivity. Observe slower movement.
			\item Press \key{Ctrl}+\key{s} -- save the existing scene. Verify the save completes.
			\item Press \key{Ctrl}+\key{c} -- print current camera parameters in snapshot format to the console.
			\item Press \key{Space} -- save a waypoint. Observe any console/UI feedback.
			\item Press \key{F1} -- switch to Pick Explore Center interaction mode.
			\item Press \key{F2} -- switch to Draw Annotation interaction mode.
			\item Press \key{F3} -- switch to Pick Annotation interaction mode.
			\item Press \key{F4} -- switch to Place Coordinate System interaction mode.
		\end{enumerate} \\
		\midrule
		\textbf{Expected Result} &
		\begin{itemize}[leftmargin=*, nosep]
			\item \key{f}: Surface texture filtering toggles visually (texture appears smoother/crisper).
			\item \key{p}: The pink explore center dot appears or disappears.
			\item \key{t}: Surfaces toggle between visible and invisible (in Comparison mode).
			\item \key{Page Up}/\key{Page Down}: Navigation speed changes noticeably.
			\item \key{Ctrl}+\key{s}: The current scene is saved without opening a dialog.
			\item \key{Ctrl}+\key{c}: Camera parameters are printed in the console output.
			\item \key{Space}: A waypoint is recorded (feedback in console or UI).
			\item \key{F1}--\key{F4}: The selected interaction mode changes in the actions drop-down.
		\end{itemize} \\
		\midrule
		\textbf{Result} & \checkbox~Pass \quad \checkbox~Fail \\
		\midrule
		\textbf{Notes} & \\[2cm]
		\bottomrule
	\end{tabularx}
	
	%----------------------------------------------------------------------------------------
	% APPENDIX: DEFECT LOG
	%----------------------------------------------------------------------------------------
	\newpage
	\section*{Defect Log}
	\label{sec:defects}
	\addcontentsline{toc}{section}{Defect Log}
	
	Use this table to record any defects found during testing.
	
	\begin{longtable}{p{1.0cm} p{1.5cm} p{3.5cm} p{5.5cm} p{1.5cm}}
		\toprule
		\textbf{ID} & \textbf{Test ID} & \textbf{Summary} & \textbf{Steps to Reproduce / Observation} & \textbf{Severity} \\
		\midrule
		\endhead
		D-01 & & & & \\[1.2cm]
		\hline
		D-02 & & & & \\[1.2cm]
		\hline
		D-03 & & & & \\[1.2cm]
		\hline
		D-04 & & & & \\[1.2cm]
		\hline
		D-05 & & & & \\[1.2cm]
		\hline
		D-06 & & & & \\[1.2cm]
		\hline
		D-07 & & & & \\[1.2cm]
		\hline
		D-08 & & & & \\[1.2cm]
		\bottomrule
	\end{longtable}
	
	\textbf{Severity levels:} Critical (blocks testing), High (major feature broken), Medium (feature partially broken), Low (minor issue / cosmetic).
	
	%----------------------------------------------------------------------------------------
\end{document}
%----------------------------------------------------------------------------------------
